\documentclass[12pt, a4paper, fontset=ubuntu]{ctexart}
\usepackage{geometry}
\usepackage{hyperref}
\usepackage{enumitem}
\usepackage{xcolor}
\usepackage{titlesec}
\usepackage{graphicx}
\usepackage{amsmath}
\usepackage{caption}
\usepackage{fontawesome5}

\geometry{left=2.5cm, right=2.5cm, top=2.5cm, bottom=2.5cm}

\hypersetup{
    colorlinks=true,
    linkcolor=black,
    urlcolor=blue
}

\titleformat{\section}{\Large\bfseries}{{\thesection}}{1em}{}
\titleformat{\subsection}{\large\bfseries}{{\thesubsection}}{1em}{}

\title{\textbf{\Huge How to start AI learning?} \\ \vspace{0.3em} \large \textit{--- for NK math \& stat}}

\author{愚者 \\ 南开大学数学科学学院2024级}
\date{2026年3月14日}


\begin{document}
\maketitle

\tableofcontents
\vfill
\noindent This article is open-sourced at \url{https://github.com/TheFool-yuzhe/nk-ai-learning-guide}.If you find it helpful,please give it a star.
\newpage


% ==================== 第一部分 ====================
\section{动机}

\subsection{不学纯数的迷茫}
笔者和很多数统人聊过，许多人处于不想做纯数、但也不知道以后做什么的迷茫困境。想做AI、但又不知道该如何去开始，同时会绕不少弯路。本文旨在从个人经历与了解信息出发，为处于迷茫、想试试入手AI的数统同学提供一些入行AI、研究生攻读AI相关的信息。碍于主笔者也只是个大二学子，许多地方尚有欠妥，如有错误还望海涵。
\subsection{关于笔者}
笔者是数院大二大类的一名学子，在大一上自知实力不足、不适合纯数，就决心找找其它方向。大一下笔者试过很多道路试错：学过python、问过计算机核心课程、打过数学建模比赛、给很多老师发过私信、用过虚拟机linux from scratch...不断试错、机缘巧合之下大一暑假进入计算机学院老师课题组，科研至今。在半年多的科研历程中，非常感谢同为南开数院本、现计算机学院研究生的师兄的指导，对如何入行AI有了很多了解。在此，笔者面向广大南开数统本科学子，给出笔者半年摸索、半年入行AI的经验与理解。

% ==================== 第二部分 ====================
\section{预备观念}

\subsection{从数学到computer science}
在接触科研之前，包括笔者在内很多数统人都有着数学人固有的一些观念和做法。比如：先学好一切预备知识，再做研究与后续。在这里笔者必须指出：当下AI发展主要是源于计算机专业及相关人员，因此许多做法、学法都是很不“数统作风”的。为了高效、少绕弯路，笔者在此先提出笔者认为数统人做AI需要先转变的观念。

\subsection{重点观念}
\begin{itemize}
    \item 数学在AI的作用远比想象中小
    \item 多实践、多敲代码、善用工具
    \item 越早入行越好
    \item 用什么学什么
    \item 进组第一
\end{itemize}

\subsection{观念解释}
\subsubsection{数学在AI的作用远比想象中小} 
一方面，AI用数学与数学系的纯数学习不同，对任一数学知识点的要求是极低的：知道有这个东西、会套公式就够用了。即使不知道，上网搜一下、AI问一下，一会也就会了够了。另一方面，当开始研究、做AI后会发现，\textbf{大多AI任务会和所学的数学知识没有任何直接关系}。换言之，当数学分析学完第十一章“多元函数微分学”、高等代数学完特征值，就已经可以入行AI了，并且越早做越好。
\subsubsection{多实践、多敲代码、善用工具}
对于数学人来说，大多敲代码很少，我们需要多敲代码练习。算法力扣使用的是C++，深度学习使用的是python。不过科研其实也不用“造轮子”式的构建，很多时候都是去GitHub上复制、AI coding（比如笔者，科研写过几百行代码了，但其实几乎0手敲）。总结下来，其一，笔者建议跟着后文推荐的“李沐深度学习”亲手敲代码练习，其二，\textbf{要善用GitHub}，其三，要善用AI。在这里笔者指出：\textbf{目前国内大模型对比国外大模型有不小差距}，如有些大模型代码能力足以胜任科研绝大部分任务（剩下的大任务上GitHub，基本都有人开源），同时有些学术网址需要VPN。不过翻墙是违规的行为，\textbf{笔者在此大力抵制这一行为}。
\subsubsection{越早入行越好}
AI对数学的知识门槛很低，并不用什么“学完数分再学”“先学好概率论、机器学习”。对于当下AI，最好的参与时间就是现在。当你学完数学分析第十一章、高等代数特征值、python基础时，你就应该开始进行AI学习尝试、联系导师进组了。AI入行需要的是遇到不会的问题再查询学习、用啥学啥、读论文，\textbf{而不是学一些数学的专业课}。
\subsubsection{用什么学什么}
现在AI的发展起步于2012AlexNet、现在的大语言模型起步于2017年的transformer，许多知识都是\textbf{这十几年提出的顶会论文上的，和纯数知识几乎毫无关系的}。因此，需要尽早学习，遇到不会顺藤摸瓜往上查询、遇到问题找开源项目，“用什么学什么”。具体来说，我需不需要先学完408再开始研究AI？不需要，用到哪个知识点搜一下知道、会用就好；我需不需要全学完概率论知识再做？毕竟AI深度学习和概率论有关系。不需要，用到哪个知识点搜一下，知道会用就好。
\subsubsection{进组第一}
AI有诸多细分领域，对某领域学习最快的方式是加入相关导师的课题组，边做边学用什么学什么。因此笔者在这里建议大家尽早进组，具体多早，大一下期中后+学个python基础（上网找网课，知道类是啥就行，这个是并行的和你数学学到哪没关系）就立马找老师。对于校内老师，笔者知道的各方面不错的是计算机学院程明明老师课题组。数学学院相关笔者打听过不少，\textbf{但从笔者个人角度出发不会选择作为日后发展方向}，故没有参与。统计学院笔者和张瑶老师聊过一次，方向很不错、交流起来人感觉也不错。其他了解过的部分方向，\textbf{不在笔者未来发展方向的考虑之内}。
\par 程明明老师主页：{\textbf{\url{https://mmcheng.net/cmm/comment-page-1/}}
程老师本人基本不带实习生，建议直接下滑联系My colleagues的小老板。有听过程老师大名的同学不要被吓到，就笔者知道的小老板而言，实习考核难度远低于研究生考核，进去实习还是不难的。Just do it。
% ==================== 第三部分 ====================
\section{学习路径推荐}

\subsection{保研与方向选择}
笔者先指出几个事情：首先，从就业赚钱角度考虑，以笔者浅薄见识的了解来看，\textbf{笔者认为硕博期间match人才计划、更容易高薪的方向，在中国几乎都在计算机学院}。其次，对于计算机相关如AILab、南京大学lamda实验室，夏令营考核机试主要分为算法和代码复现。算法需要学习前置数据结构、刷力扣、语言为C++，而深度学习平常使用都是python，同时和算法没有任何关系。而进组实验的过程中少不了复现GitHub代码、读论文、分析论文原理复现。故笔者建议，一方面先学python（知道啥是类）然后进组，另一方面学习数据结构、C++、刷算法题（这个不着急，因为笔者在目前半年多科研经历中没用到一次算法相关，主要是应对夏令营考核，需要自行安排时间）。
\par 由于笔者也仅对科研相关的cv(计算机视觉)、nlp(自然语言处理)领域有一点了解，故接下来推荐的学习路径主要是和cv、nlp强相关。不过一方面，cv、nlp是AI、深度学习火起来的开端；另一方面，程老师课题组方向是cv，张瑶老师课题组方向是nlp。
\subsection{线性时间规划}
1、学习python，知道类是啥就行（不需要学torch）。简单了解Linux系统。
\par 2、跟着“动手学深度学习（李沐）”学习当下AI的深度学习原理、代码原理。在这期间他会逐步教你torch、AI的实现原理。\url{https://zh-v2.d2l.ai/}
\begin{center}
\begin{minipage}{0.12\textwidth}
    \centering
    \includegraphics[width=\textwidth]{AIlearning/1.png}
    \captionof{figure}{李沐深度学习}
\end{minipage}
\end{center}
\par 3、如果跟cv/nlp课题组，此时应该投简历进组。
\par 4、并行的推荐关注以下博主精读论文，同时如对CNN、attention、ViT不了解原理，推荐去\textbf{哔哩哔哩搜索讲解视频}：跟李沐学AI、3Blue1Brown。平常有空\textbf{多刷刷AI相关内容}，看看现在AI啥发展。
\begin{center}
\begin{minipage}{0.2\textwidth}
    \centering
    
\includegraphics[width=\textwidth]{AIlearning/2.png}
\end{minipage}
\end{center}
\begin{center}
\begin{minipage}{0.2\textwidth}
    \centering
    
\includegraphics[width=\textwidth]{AIlearning/3.png}
\end{minipage}
\end{center}
\par 5、从笔者浅薄见识个人见解，不推荐学数学建模比赛，对升学、AI学习作用很小，时间投入产出比很低。平常多搜搜AI相关知识，多去GitHub逛逛。
\par 6、推荐使用ZOTERO进行论文管理，以及arxiv进行论文搜索，善于用知乎、CSDN寻找问题解决答案。
\\ZOTERO:\url{https://zhuanlan.zhihu.com/p/669287007}
\\\url{https://www.zotero.org/download/}
\\必备扩展：Zotero Connector
\\ arxiv:\url{https://arxiv.org/search/cs}
% ==================== 第四部分 ====================
\section{常见问题QA}
\subsection{进组是否太早？啥也不会怎么参与？}
A：笔者进组的时候更菜，深度学习一点不会进组后学的，相当于就会个Linux和python到类（torch都不会）。跟着边做边学就好了，重要的是抱有一颗持续学习的心。
\subsection{被老师拒绝、无法进组怎么办？}
A：简历写好点、多发点简历就好了，笔者也好几次被老师已读不回过。此外校内老师还是有几个会对数统学子宽容度更高一些，觉得学数学的聪明的。
\subsection{和课内专业课时间平衡怎么办？}
A:自我权衡了，从笔者角度当然更希望大家都多尝试一些非专业课的东西“一群人努力提高分数线"，大家都学点别的=大家排名不变，数院绩点排名太吓人了。此外程明明老师课题组去向综合考虑是很强的（以及氛围、就读体验），很可能卷一大顿专业课拼外保，最后研究生毕业发现不如在程老师课题组。程老师大组不看重rk排名只看实力，有保研名额就行。
\subsection{感觉要学的好多压力好大。}
A:路虽远，行则将至。一方面，上面学习所需时间并不长。快的话两周学python、五周学深度学习、一周联系导师进组，两个月就有入门基础。剩下时间简单过一下数据结构，刷刷力扣，跟着课题组多敲敲代码、读读论文，半年就能有不小进步、有去不错去向的可能。另一方面，相信自己考到南开数统的智商。AI相关理解成本不高。当时深度学习那本书笔者只用一周就简单过了一遍，相比于数学专业课难度低很多。
\subsection{那数学专业相关学这么多是否没用？}
A:并非，从笔者接触科研来说：一方面，笔者做computer science时往往要一心多用、多窗口任务。在短视频和vibe coding的当下，这更加影响专注力。数学专业课的学习培养笔者心无杂念、一心一用专心研究一个问题的能力，毕竟三心二意数学题真的做不出来（对于笔者）。这份attention笔者认为在科研中和其他地方都是很宝贵的。另一方面，国内本科教育问题很多，就从事AI行业为导向，其实别的专业哪怕是计算机专业，课程设置导致无用内容学习、时间浪费也不少，相对来说数统课还是少。而且当下AI最火的华人，21世纪单篇引用量最高论文作者，何恺明，本科也是学理科出身而非计算机科班出身（清华物理系）。
\subsection{学数学专业课程/408毕竟还是有用，要不还是先学概率论什么的？}
A:这是个\textbf{时间性价比}的问题。一天只有24小时，先学这些别的，会导致最后学AI、接触科研时间就少了。从笔者经验来看，这是捡了芝麻丢了西瓜、得不偿失。
% ==================== 第五部分 ====================
\section{结语}
以上是笔者半年摸索、半年AI探路的心得，旨在为还在迷茫中的数统学子指一条参与AI、除了专业课外有努力地方的道路。本文只是个起步建议，至于进组后怎么办，其一作者也只是个大二学子，其二相信大家也会有自己的感悟与方向、无需笔者赘述，其三进组之后\textbf{多和研究生师兄师姐们聊聊}，大家会豁然开朗，得到远比笔者提供的更多、更有用的信息。最后总结一下最重要的三点：
\begin{itemize}
    \item \textbf{早开始、早进组}
    \item \textbf{用啥学啥}
    \item \textbf{Just do it}
\end{itemize}
\par 虽然笔者很容易盒出来，但不推荐大家飞书私信问笔者问题。\textbf{笔者平常学业繁忙}没空回答、不会回复关于本文问题。和笔者比较熟的可以直接微信私聊笔者，学长们有相关经历经验欢迎指正，笔者收集收集以后或许会出V2、V3。金融等其它方面也期待有大神开源讲解，为大家指路～
% ==================== 第六部分 ====================
\section{附录}
\subsection{关于附录}
笔者后续向几位已保研/出国的学长请教，得到了更全面、中肯的建议，在此作为补充。特别感谢为本文提供宝贵建议的学长们。（别问为什么没提学姐，数院男女比……）
\subsection{关于大模型}
笔者和学长鼓励大家去\textbf{详细搜索国外大模型如ChatGPT、Claude的效果}，仔细了解目前\textbf{国内大模型对比美国的差距}，比如关注一下\textbf{ChatGPT写出来代码效果、信息问题检索能力}。目前大家\textbf{熟知}的大模型中，有一些大模型已能从规划学习到具体实施，抹平很多信息差以及作为高效生产力，比如按照需求生成高质量代码。如果大家平常用的大模型做不到这么好用，可以思考一下是否要\textbf{换一个大模型使用}。不过翻墙挂梯子是错误的行为，在这里笔者\textbf{坚决抵制翻墙}。
\subsection{关于强化学习、AI学习和北美直博}
\noindent{\small \textit{by NON\enspace|\enspace 2021NKmath $\Rightarrow$ UCB IEOR PhD}}
\subsubsection{关于本文补充}
\textit{1. 对于ai其中的不少领域，不仅仅是靠数分高代，stochastic和statistics的相关知识也是必要的，当然尽早接触项目是好事，但是不要忽视这些知识。}
\par \textit{2. 可以再多强调一下ai的使用，从规划学习到具体实施，ai能抹平很多信息差以及作为高效生产力。}
\subsubsection{关于强化学习入门}
\textit{RL部分的话，推荐的是《强化学习的数学原理》，b站上就有视频，也有GitHub。}
\subsubsection{关于AI其他领域入门}
\textit{领域蛮多的，偏优化、系统、安全、RL、理论等等，感觉大家的问题很多在于不了解这个领域以及不知道怎么入门？ai能给出比较好的一些学习建议吧，然后就可以去联系相关领域的老师去尝试做一下科研这样}
\par 笔者：这位学长实力非常强悍，大四2.5篇A会一作申的伯克利phd（国内保研应该哪都随便去了）。由于\textbf{时间精力}、个人水平，以及综合考虑\textbf{升学}的前提下，笔者建议优先考虑校内有导师+个人能接受的方向，简单了解后进组学习
\subsubsection{关于出国}
1、不同于国内学硕/直博低区分度、甚至直博bar更低。国外PhD直博是显著比硕士难申请的，但同时PhD其一0花销、有工资cover平常生活（硕士很贵），其二认可度高于硕士不少。并且硕士/PhD申请看重的不太相同，下面分享是以申请\textbf{泛AI、泛cs北美PhD}为主。
\par 2、重要程度\textbf{推荐信}>paper $\gg$ GPA>SoP>语言，后三过关不难看就行。“\textbf{connection is all you need.}”推荐信是最重要的、尤其需要强推＆最好是大牛/有私交。推荐信的一个重要来源是跟着导师做科研，故学长推荐流程是早进组，“校内导师$\Rightarrow$外校导师$\Rightarrow$美国导师” 流程，三段科研三段产出三封推荐信。
\par 3、有时间不建议找中介，多做科研和个人科研水平是重点。
\subsection{关于保研、夏令营}
\noindent{\small \textit{by Qiu\enspace|\enspace 2022NKstat $\Rightarrow$ AILAB $\times$ FDU PhD}}
\subsubsection{关于科研-专业课时间平衡}
\textit{“不过你是学分绩比较高，所以可以早点进组，有的学分绩不高的，如果提早进组得考虑学分绩的问题”}
\par 笔者：rk排名确实是个问题，不过进组接触AI确实收益也不小，或许可以考虑像笔者一样扔一些选修课去大四，毕竟进组收益很高。
\subsubsection{关于机试考核与硕博选择}
\textit{夏令营的话，就是可以选《算法基础》那个课，基本上基础部分就没问题了，不过还是要\textbf{早点想好读博还是读硕}吧，这两个准备的重点不太一样。}
\par 补充：这个有前文的。国内硕士在考察时更重机试、上机代码考察、rk排名，直博更重科研经历与产出、rk排名、和导师提前打招呼。同时硕博培养方式又不太一样，硕士更灵活压力更小，直博有论文、延毕压力但五年拿到博士学位上限更高，而且不少院系直博门槛更低、学硕名额少甚至没有。提前想清楚硕博很重要，\textbf{不要一昧追title},尤其对于博士，\textbf{导师极大决定未来几年生活幸福度和就业发展。}
\subsubsection{关于比赛、夏令营笔试与面试}
1、对于AI、偏cs系保研，除ACM、CMC国一外基本上用处很小，从\textbf{时间投入产出比}角度考量，笔者认为将有限时间用到别处可能会更好。
\par 2、算法对应的是笔试，AI知识和代码复现、读论文能力对应的是面试。\textit{“面试的考查范围就方差很大了，有的学校喜欢考一些基础的科班的知识，有的喜欢问最新的技术，这也看面试的老师的兴趣。”}
\subsection{笔者总结}
笔者总结一下上文关于升学部分：保研硕士看重\textbf{机试、rk排名、其他}，保研直博看重\textbf{科研经历与产出、rk排名、其他}，美国直博看重\textbf{推荐信}（科研经历带来）、\textbf{科研产出}。进组科研时间会占用专业课学习时间，可能会导致rk排名问题。但硕博选择同时也很重要，而是否适合博士和科研又需要通过进组科研检验，并且进组能学到很多东西。因为延毕压力、不适合科研等，每天很焦虑甚至抑郁的博士生有不少，不能为了title盲追直博。故笔者在这里建议：大二时少选一些选修课，早日进组尝试科研以学习和检验自己是否适合科研，同时兼顾专业课学习情况进行科研上的时间分配。在科研经历中想好未来发展方向，硕/博，是否出国，再做下一步抉择。如果选择读博，导师选择非常重要。\textbf{导师极大决定未来几年生活幸福度和就业发展。}
\bigskip
\bigskip
\bigskip
\begin{center}
    {\large \textit{祝大家都能找到自己未来的方向，走上自己喜欢的道路，活成自己期待的样子，开心过好人生每一天！}}
\end{center}
\end{document}

